%! Catch-Up Packet — Lessons 1.0–2.2 (for a student joining mid-course)
% Standalone REVIEW handout for a new (exchange) student joining during Lesson 2.2: an abridged
% walk through every lesson so far, in plain English with short sentences. Read-only, not turned
% in, not scored, and with NO key (user decision 2026-10-22) — every worked example is printed
% in full with ordinary \[ \] displays, and there is no \namedateperiod.
% Authored on the unit study-guide model (unit02/review/study_guide: 10pt, \parthead strips);
% it is NOT a lesson component and is merged into no packet.
\documentclass[10pt]{article}
\usepackage{saar-article}
\usepackage{saar-boxes}

% Part-divider strip (same macro the unit tests define)
\newcommand{\parthead}[1]{%
  \vspace{6pt}\noindent
  \begin{headlinebox}{cerulean}{\color{white}\bfseries #1}\end{headlinebox}%
  \vspace{4pt}}

\begin{document}

\pageheader{Catch-Up Packet}{Lessons 1.0--2.2 --- Study Design and Describing One Variable}

\vspace{0.10in}
\begin{remindbox}
\small
\textbf{Welcome to the class!} This packet shows you what we learned before you arrived. It is
short. It covers \textbf{Unit 1} (how to collect data, Lessons 1.0--1.8) and the start of
\textbf{Unit 2} (how to describe data, Lessons 2.1--2.2).
\textbf{It is not homework.} You do not turn it in, and it is not graded.
\textbf{How to use it:} (1) Read the word list in Part 1. Cover the meanings and test yourself.
(2) Read each example in Part 2 slowly, one line at a time. (3) Use the checklist at the end.
Mark anything you do not understand, and ask your teacher.
\end{remindbox}

% ======================================================================
\parthead{Part 1 --- The Words You Need \hfill Lessons 1.0--2.2}

\small
Each meaning uses short, simple sentences. Learn the word on the left --- we use it in every class.

\vspace{3pt}
\renewcommand{\arraystretch}{1.15}
% Three tables, not one: a tabularx cannot break across a page. Split at lesson boundaries.
\boxguard[14]
\begin{tabularx}{\linewidth}{@{} c >{\raggedright\arraybackslash\bfseries}p{3.4cm} X @{}}
\rowcolor{frost}
\textbf{Lsn} & \textbf{Word} & \textbf{What it means} \\
\midrule
1.0 & Individual & One person or thing we collect data about. One \emph{row} of a table. \\
1.0 & Variable & One thing we record about each individual. One \emph{column} of a table. \\
1.0 & Categorical variable & The answer is a \emph{label} or a group name. Example: favorite
      sport. \\
1.0 & Quantitative variable & The answer is an \emph{amount} with units. Example: hours of
      sleep. \\
1.1 & Data cycle & Every study has four steps, in order: \textbf{ask} a question,
      \textbf{collect} data, \textbf{organize} it, \textbf{analyze} it and say what it means. \\
1.1 & Statistical question & A question with answers that \emph{vary}. It names the group and
      what we will record. \\
1.1 & Relative frequency & A count divided by the total. Write it as a decimal or a percent. \\
1.2 & Population & \emph{Everyone} the question is about. \\
1.2 & Sample & The part of the population we \emph{actually} collect data from. \\
1.2 & Parameter / statistic & A number about the \emph{population} (usually unknown) / a
      number from the \emph{sample} (we compute it). \\
1.2 & Sampling variability & Two good samples give different numbers. This is normal. It is
      not a mistake. \\
\end{tabularx}

\vspace{4pt}
\boxguard[16]
\begin{tabularx}{\linewidth}{@{} c >{\raggedright\arraybackslash\bfseries}p{3.4cm} X @{}}
\rowcolor{frost}
\textbf{Lsn} & \textbf{Word} & \textbf{What it means} \\
\midrule
1.3 & Random sample & \emph{Chance} chooses who is in the sample, not a person. \\
1.3 & Convenience sample & We ask whoever is easy to reach. This is \emph{not} random. \\
1.4 & Bias & The method leans the same way every time. It always guesses too high, or always
      too low. \\
1.5 & Observational study & We watch or ask. We do \emph{not} put anyone in a group. \\
1.5 & Explanatory / response variable & The variable we think causes a change / the variable
      we measure to see the change. \\
1.5 & Association & Two variables change together. This does \emph{not} prove that one causes
      the other. \\
1.5 & Lurking variable & A hidden third variable that affects both variables. \\
1.6 & Experiment & The researcher \emph{assigns} the groups and gives a treatment. Only an
      experiment can show \emph{cause}. \\
1.6 & Treatment / control group & The group that gets the treatment / the group that does not.
      We compare them. \\
1.6 & Placebo & A fake treatment that looks real. \\
1.6 & Blinding & People do not know who is in which group. \textbf{Double-blind:} the
      subjects \emph{and} the measurers do not know. \\
1.7 & Collection method & \emph{How} we get the data: survey, interview, focus group,
      observation, or existing records. \\
1.8 & Survey instrument & The exact questions we ask, in order, with the answer choices. \\
\end{tabularx}

\vspace{4pt}
\boxguard[16]
\begin{tabularx}{\linewidth}{@{} c >{\raggedright\arraybackslash\bfseries}p{3.4cm} X @{}}
\rowcolor{frost}
\textbf{Lsn} & \textbf{Word} & \textbf{What it means} \\
\midrule
2.1 & Distribution & The values a variable takes, and how often it takes each one. \\
2.1 & Dot plot / stemplot / histogram & Three pictures of quantitative data. A dot plot has one
      dot per individual. A stemplot splits each number into a stem and a leaf. A histogram
      has bars. \\
2.1 & Skewed right / left & One side has a long \emph{tail}. The tail gives the name: a tail
      to the right is skewed \emph{right}. \\
2.1 & Outlier & A value far away from the other values. \\
2.2 & Mean ($\bar{x}$) & Add all the values. Divide by how many there are. \\
2.2 & Median & The middle value when the data are in order. \\
2.2 & Range / IQR & Maximum $-$ minimum / $Q_3 - Q_1$, the spread of the middle half. \\
2.2 & Standard deviation ($s$) & About how far a typical value is from the mean. It measures
      \emph{spread}, not size. \\
\end{tabularx}

\vspace{6pt}
\boxguard[14]
\begin{blurbbox}[Words in our questions]{goldacc}
\small
\textbf{Describe} --- tell what you see. \quad
\textbf{Interpret} --- say what a number means in real life. \quad
\textbf{Justify} --- give the reason (``because \dots''). \quad
\textbf{In context} --- use the real words and units, for example ``$13$ minutes for these
students,'' not just ``$13$.''
\end{blurbbox}
\normalsize

% ======================================================================
\vspace{4pt}
\boxguard[22]
\parthead{Part 2 --- The Ten Skills, Worked \hfill Unit 1, then Unit 2}

\small
Each box shows one skill, with all the work. Each box ends with \textbf{the trap} --- the most
common mistake. Every example is about Maple Grove High School, which has $1{,}200$ students.

\vspace{4pt}
\boxguard[14]
\begin{notesbox}{Skill 1 --- Is a variable categorical or quantitative? (Lesson 1.0)}
\small
Ask one question: \emph{does math with this variable make sense?} If yes, it is quantitative.
If no, it is categorical.

\vspace{2pt}
\renewcommand{\arraystretch}{1.2}
\begin{tabularx}{\linewidth}{@{} X c X @{}}
\rowcolor{frost}
\textbf{Variable} & \textbf{Type} & \textbf{Why} \\
\midrule
Minutes to get to school & Quantitative & An average of minutes makes sense. \\
Bus number & Categorical & The number is a name. ``Average bus $7.5$'' means nothing. \\
Favorite subject & Categorical & It is a label. We count it. \\
\end{tabularx}

\vspace{3pt}
\textbf{The trap:} ``It has digits, so it is quantitative.'' Bus numbers, phone numbers, and
ZIP codes have digits, but they are \emph{labels}.
\end{notesbox}

\vspace{3pt}
\boxguard[14]
\begin{notesbox}{Skill 2 --- Ask a statistical question, and find a percent (Lesson 1.1)}
\small
``How many minutes does Ana sleep?'' has \emph{one} answer, so it is not statistical.
\emph{``How many minutes do Maple Grove students sleep on a school night?''} names a group, and
the answers vary. That is a \textbf{statistical question}.

\vspace{2pt}
In a class of $25$, $10$ students walk to school. The relative frequency is
\[ \frac{10}{25} = 0.40 = 40\% \]
\textbf{The trap:} a question with no group, or a question with only one answer.
\end{notesbox}

\vspace{3pt}
\boxguard[8]
\begin{notesbox}{Skill 3 --- Population or sample? Parameter or statistic? (Lesson 1.2)}
\small
The principal picks $60$ of the $1{,}200$ students at random. $18$ of the $60$ walk to school.

\vspace{2pt}
\textbf{Population:} all $1{,}200$ students. \quad \textbf{Sample:} the $60$ students asked.
\[ \frac{18}{60} = 0.30 = 30\% \qquad\text{so we estimate}\qquad 0.30 \times 1200 = 360
   \text{ students walk} \]
The $30\%$ is a \textbf{statistic}: it comes from the sample. The true percent for all
$1{,}200$ students is the \textbf{parameter}. We do not know it --- the $360$ is an
\textbf{estimate}.

\vspace{2pt}
\textbf{The trap:} a second random sample gets $35\%$, and someone says the first survey was
wrong. It was not. Different samples give different numbers. That is \textbf{sampling
variability}.
\end{notesbox}

\vspace{3pt}
\boxguard[16]
\begin{notesbox}{Skill 4 --- Choose a sample with chance (Lesson 1.3)}
\small
\renewcommand{\arraystretch}{1.2}
\begin{tabularx}{\linewidth}{@{} >{\bfseries}p{2.6cm} X @{}}
\rowcolor{frost}
\textbf{Method} & \textbf{How it works} \\
\midrule
Simple random & Give every student a number. Draw numbers at random. \\
Systematic & Pick a random start on the list. Then take every $k$th student. \\
Stratified & Split into groups that are \emph{alike inside} (like grades). Take some at random
  from \emph{every} group. \\
Cluster & Split into groups that are each a \emph{mix} (like advisory groups). Pick a few whole
  groups at random. Ask everyone in them. \\
\end{tabularx}

\vspace{3pt}
\textbf{Systematic:} we want $60$ of $1{,}200$, so $k = 1200 \div 60 = 20$. Start at a random
number, like $7$, and take students $7, 27, 47, 67, \dots$

\vspace{2pt}
\textbf{Stratified:} take the same fraction, $\frac{60}{1200} = \frac{1}{20}$, of every grade.

\vspace{-2pt}
\begin{tabularx}{\linewidth}{@{} X c c c c | c @{}}
\rowcolor{frost}
\textbf{Grade} & \textbf{9} & \textbf{10} & \textbf{11} & \textbf{12} & \textbf{Total} \\
\midrule
Students in the grade & $320$ & $300$ & $300$ & $280$ & $1{,}200$ \\
Number in the sample ($\div 20$) & $16$ & $15$ & $15$ & $14$ & $60$ \\
\end{tabularx}

\vspace{3pt}
\textbf{The trap:} asking whoever walks by is \emph{not} random. That is a
\textbf{convenience sample}. Chance must choose.
\end{notesbox}

\vspace{3pt}
\boxguard[16]
\begin{notesbox}{Skill 5 --- Name the bias (Lesson 1.4)}
\small
Ask: \emph{could this person be chosen, and did they answer honestly?}

\vspace{2pt}
\renewcommand{\arraystretch}{1.2}
\begin{tabularx}{\linewidth}{@{} >{\bfseries}p{2.6cm} X X @{}}
\rowcolor{frost}
\textbf{Type of bias} & \textbf{What went wrong} & \textbf{Example} \\
\midrule
Undercoverage & Some people were never on the list. & An online survey misses students with
  no internet. \\
Nonresponse & People were chosen but did not answer. & Only $20$ of $100$ emails come back. \\
Response bias & People answered, but the question pushed them. & ``Don't you agree lunch is
  too short?'' (a \emph{leading} question) \\
\end{tabularx}

\vspace{3pt}
\textbf{The trap:} ``Ask more people and the bias goes away.'' No. A bigger sample with the
same bad method makes the same mistake. Only a \emph{better method} removes bias.
\end{notesbox}

\vspace{3pt}
\boxguard[16]
\begin{notesbox}{Skill 6 --- Association or cause? (Lessons 1.5 and 1.7)}
\small
\textbf{Ask: who made the groups?} If the people chose (or the groups already existed), it is
an \textbf{observational study}. If the researcher assigned the groups at random, it is an
\textbf{experiment}.

\vspace{2pt}
A survey finds that Maple Grove students who play sports have higher grades. The students chose
to play, so this is an observational study. We may say sports and grades are
\textbf{associated}. We may \emph{not} say sports \textbf{cause} higher grades. A
\textbf{lurking variable} could explain it: for example, students with more free time may do
both.

\vspace{2pt}
\textbf{The trap:} ``The groups are different, so one thing caused the other.'' Only an
experiment can show cause.
\end{notesbox}

\vspace{3pt}
\boxguard[18]
\begin{notesbox}{Skill 7 --- Plan an experiment (Lesson 1.6)}
\small
Does a text-message reminder help students turn in homework? A computer randomly puts $40$
students in a \textbf{treatment group} (they get the text) and $40$ in a \textbf{control
group} (no text).
\[ \text{text: } \frac{24}{40} = 60\% \qquad \text{no text: } \frac{12}{40} = 30\% \qquad
   60 - 30 = 30 \text{ percentage points} \]
Chance made the groups, so the groups started out alike. So we \emph{may} say the text
\textbf{caused} the increase.

\vspace{2pt}
\renewcommand{\arraystretch}{1.15}
\begin{tabularx}{\linewidth}{@{} >{\bfseries}p{2.6cm} X @{}}
\rowcolor{frost}
\textbf{Principle} & \textbf{In this experiment} \\
\midrule
Comparison & Two groups: text and no text. \\
Randomization & A computer, not the teacher, picks the groups. \\
Replication & $40$ students in each group, not just a few. \\
Control & Same homework, same week, same teacher. Only the text is different. \\
\end{tabularx}

\vspace{3pt}
\textbf{More words:} a \textbf{placebo} is a fake treatment. \textbf{Blinding} means people
do not know their group. A \textbf{block} is a group of similar subjects (for example, one
grade); we randomize \emph{inside} each block.

\vspace{2pt}
\textbf{The trap:} ``The control group is wasted.'' No --- without it, there is nothing to
compare to.
\end{notesbox}

\vspace{3pt}
\boxguard[16]
\begin{notesbox}{Skill 8 --- Write good survey questions (Lessons 1.7 and 1.8)}
\small
\renewcommand{\arraystretch}{1.2}
\begin{tabularx}{\linewidth}{@{} >{\bfseries}p{3.0cm} X X @{}}
\rowcolor{frost}
\textbf{Rule} & \textbf{Bad question} & \textbf{Better question} \\
\midrule
One idea per question & ``Is the cafeteria clean and cheap?'' & Ask two questions: one about
  clean, one about cheap. \\
Neutral words & ``Don't you agree the library is great?'' & ``How would you rate the
  library?'' \\
Choices fit everyone & ``$0$--$2$ hours'' and ``$2$--$4$ hours'' (where does $2$ go?) &
  ``$0$--$1$,'' ``$2$--$3$,'' ``$4$ or more'' \\
\end{tabularx}

\vspace{3pt}
\textbf{Pick the method:} numbers from many people $\rightarrow$ \textbf{survey}; reasons in
their own words $\rightarrow$ \textbf{interview}; what people really do $\rightarrow$
\textbf{observation}. \textbf{The trap:} good questions do not fix a bad sample. You still need chance to choose who
answers.
\end{notesbox}

\vspace{3pt}
\boxguard[26]
\begin{notesbox}{Skill 9 --- Read a display and describe the distribution (Lesson 2.1)}
\small
Fifteen students recorded how many minutes it takes them to get to school. Here are the same
$15$ values two ways.

\vspace{3pt}
\begin{minipage}[t]{0.30\linewidth}
\vspace{0pt}
\textbf{Stemplot}

\vspace{2pt}
\renewcommand{\arraystretch}{1.1}
\begin{tabular}{r|l}
0 & 5\;7\;8 \\
1 & 0\;0\;2\;2\;3\;5\;5\;8 \\
2 & 0\;2\;5 \\
3 & \\
4 & 0 \\
\end{tabular}

\vspace{3pt}
\textbf{Key:} $1\,|\,3$ means $13$ minutes
\end{minipage}\hfill
\begin{minipage}[t]{0.66\linewidth}
\vspace{0pt}
\textbf{Dot plot}

\vspace{2pt}
\begin{tikzpicture}[x=0.24cm, y=0.22cm]
  \draw[linegray] (0,0) -- (42,0);
  \foreach \t in {0,5,...,40} {
    \draw[linegray] (\t,-0.3) -- (\t,0.3);
    \node[below, font=\scriptsize] at (\t,-0.3) {\t};}
  \node[below, font=\scriptsize] at (21,-2.0) {Minutes to get to school};
  \foreach \v/\h in {5/1, 7/1, 8/1, 10/1, 10/2, 12/1, 12/2, 13/1, 15/1, 15/2, 18/1,
                     20/1, 22/1, 25/1, 40/1}
    \fill[cerulean] (\v,\h) circle (0.35);
\end{tikzpicture}
\end{minipage}

\vspace{4pt}
\textbf{Describe it with four things:} \textbf{shape}, anything \textbf{unusual}, the
\textbf{center}, and the \textbf{spread} --- in context. There are $15$ values, so the median
is the $8$th value in order: $13$ minutes. The range is $40 - 5 = 35$ minutes.

\vspace{2pt}
\emph{The travel times for these $15$ students are skewed right. There is a gap from $25$ to
$40$ minutes, and $40$ minutes is a possible outlier. The median is $13$ minutes, and the
times go from $5$ to $40$ minutes.}

\vspace{3pt}
\textbf{The trap:} naming the skew by where the dots pile up. The pile is on the left, but the
long \emph{tail} goes to the right. \textbf{The tail names the skew}, so it is skewed right.
\end{notesbox}

\vspace{3pt}
\boxguard[22]
\begin{notesbox}{Skill 10 --- Center and spread: mean, median, IQR, standard deviation (Lesson 2.2)}
\small
Seven students' hours of homework last week, in order: \quad
$3$, $5$, $6$, $7$, $8$, $9$, $11$.

\vspace{2pt}
\textbf{Center.}\vspace{-4pt}
\[ \bar{x} = \frac{3 + 5 + 6 + 7 + 8 + 9 + 11}{7} = \frac{49}{7} = 7 \text{ hours}
   \qquad \text{median} = 7 \text{ hours (the 4th value)} \]
\textbf{Quartiles.} Leave the median out. $Q_1$ is the middle of $3, 5, 6$, and $Q_3$ is the
middle of $8, 9, 11$:
\[ Q_1 = 5 \qquad Q_3 = 9 \qquad \text{IQR} = 9 - 5 = 4 \text{ hours} \qquad
   \text{range} = 11 - 3 = 8 \text{ hours} \]
\textbf{Five-number summary:} $3$, $5$, $7$, $9$, $11$ (min, $Q_1$, median, $Q_3$, max).
\textbf{Standard deviation.} Subtract the mean from each value, square, add, divide by
$n - 1$, then take the square root.

\vspace{2pt}
\renewcommand{\arraystretch}{1.2}
\begin{tabularx}{\linewidth}{@{} >{\bfseries}l *{7}{>{\centering\arraybackslash}X} | c @{}}
\rowcolor{frost}
 & & & & & & & & \textbf{Sum} \\
Value $x$ & $3$ & $5$ & $6$ & $7$ & $8$ & $9$ & $11$ & $49$ \\
Deviation $x - \bar{x}$ & $-4$ & $-2$ & $-1$ & $0$ & $1$ & $2$ & $4$ & $0$ \\
Squared $(x - \bar{x})^2$ & $16$ & $4$ & $1$ & $0$ & $1$ & $4$ & $16$ & $42$ \\
\end{tabularx}
\[ s^2 = \frac{42}{7 - 1} = 7 \qquad s = \sqrt{7} \approx 2.6 \text{ hours} \]
\textbf{Say it in context:} \emph{A typical student's homework time is about $2.6$ hours away
from the mean of $7$ hours.}

\vspace{3pt}
\textbf{The trap:} ``A bigger standard deviation means bigger values.'' No. The standard
deviation measures \textbf{spread}, not size. If every student did $10$ more hours, the mean
would go up by $10$, but $s$ would stay $2.6$.
\end{notesbox}
\normalsize

\vspace{4pt}
\boxguard[5]
\begin{remindbox}
\small
\textbf{You are caught up when you can do Skills 1--10 without looking.} Not sure about one?
Circle it and ask your teacher.
\end{remindbox}

\end{document}
